\usecasetable
{UC-05}
{Create Booking}
{Customer, Owner}
{Permette ad un Customer di prenotare posti, servizi extra e parcheggi presso una spiaggia in una certa data. Permette inoltre ad un Owner di registrare manualmente una prenotazione per un cliente telefonico.}
{L'utente deve aver effettuato il login. Se l'utente è un Customer, non deve essere bannato dalla spiaggia selezionata. Se l'utente è un Owner, deve possedere una spiaggia attiva. La spiaggia selezionata (o posseduta dall'Owner) deve essere nello stato ATTIVO.}
{Viene creato un nuovo record di prenotazione nel database con un prezzo totale calcolato e uno stato iniziale (PENDING per i Customer, CONFIRMED per gli Owner).}
{
	\begin{enumerate}[nosep, leftmargin=*]
		\item Il Customer seleziona una spiaggia (partendo da UC-04) e sceglie una data;
		\item Il sistema verifica che il Customer non sia bannato da quella spiaggia;
		\item Il sistema mostra la mappa o la lista delle disponibilità per la data scelta;
		\item Il Customer seleziona i posti spiaggia desiderati (es. ombrelloni, lettini);
		\item Il Customer seleziona eventuali extra (lettini, sdraio, sedie);
		\item Il Customer inserisce il numero di parcheggi desiderati;
		\item Il sistema calcola e mostra il prezzo totale riepilogativo;
		\item Il Customer conferma la prenotazione;
		\item Il sistema salva la prenotazione e le assegna lo stato "PENDING", mostrando un messaggio di successo.
	\end{enumerate}
}
{
	\begin{enumerate}[nosep, leftmargin=36px]
		\item[\textbf{1a.}] \textbf{Flusso Owner (Prenotazione Telefonica):}
		\begin{enumerate}[nosep, leftmargin=*]
			\item[i.] L'Owner accede al pannello gestionale della propria spiaggia e avvia una nuova prenotazione;
			\item[ii.] L'Owner seleziona la data, i posti, gli extra e i parcheggi (come nei passi 3-6 del Main Flow);
			\item[iii.] L'Owner inserisce obbligatoriamente il Nome e il Numero di Telefono del cliente chiamante;
			\item[iv.] Il sistema calcola e mostra il totale;
			\item[v.] L'Owner conferma la prenotazione;
			\item[vi.] Il sistema salva la prenotazione assegnandole automaticamente lo stato "CONFIRMED".
		\end{enumerate}
		\item[\textbf{2a-e.}] \textbf{Utente Bannato:} Il sistema rileva che il Customer è nella blacklist della spiaggia e blocca l'operazione mostrando il messaggio: "Non sei autorizzato a prenotare in questa struttura.".
		\item[\textbf{4/6a-e.}] \textbf{Disponibilità Esaurita:} Se durante la selezione i posti o i parcheggi scelti vengono esauriti (es. prenotati da un altro utente nello stesso istante), il sistema avvisa l'utente e lo invita a modificare la scelta.
	\end{enumerate}
}
{
	\begin{enumerate}[nosep, leftmargin=*]
		\item[•] Le prenotazioni dei Customer partono sempre dallo stato PENDING e richiedono un passaggio a CONFIRMED da parte dell'Owner.
		\item[•] Le prenotazioni create dall'Owner sono considerate già confermate (CONFIRMED). 
		\item[•] L'Owner non può selezionare la spiaggia: il sistema forza la prenotazione solo sullo stabilimento di sua proprietà.
		\item[•] Il calcolo del totale si basa sul tariffario (Season/Zone) impostato per quella data.
	\end{enumerate}
	
	\begin{enumerate}[leftmargin=*]
		\item[•] \textbf{Mockup}: si veda \autoref{fig:mockup3} - Flusso di Prenotazione
		
		\item[•] \textbf{Testing}: le classi che si occupano dei test sono \hyperref[itm:application_booking]{\texttt{CreateBookingServiceTest}} e \hyperref[itm:application_booking]{\texttt{CreateManualBookingServiceTest}} nel package \texttt{service}; \hyperref[itm:domain_booking_pricecalculator]{\texttt{BookingTest}}, \hyperref[itm:domain_booking_pricecalculator]{\texttt{BookingParkingTest}} e \hyperref[itm:domain_booking_pricecalculator]{\texttt{PriceCalculatorTest}} nel package \texttt{domain}.
	\end{enumerate}
}